\begin{frame}[t]{SVAR-модель Sims--Zha}
\small

Sims--Zha используют \textbf{структурную VAR-модель}: приведённая VAR описывает динамику данных, а структурная форма нужна для интерпретации шоков.

\vspace{0.35em}

\begin{columns}[T,onlytextwidth]

\begin{column}{0.48\textwidth}
\textbf{Вектор переменных}

\vspace{0.3em}

\begin{itemize}
\setlength\itemsep{0.32em}
\item \textbf{Денежный блок:} $R$ — краткосрочная ставка; $M$ — M2 или совокупные резервы.
\item \textbf{Реальная экономика:} $y$ — выпуск; $W$ — заработная плата.
\item \textbf{Цены:} $P_y$ — дефлятор выпуска; $P_{im}$ — цены промежуточных материалов.
\item \textbf{Информационные переменные:} $Tbk$ — банкротства; $P_{cm}$ — сырьевые цены.
\end{itemize}

\vspace{0.25em}

\footnotesize
Идентификация задаётся ограничениями на $\Gamma_0$.
\end{column}

\begin{column}{0.48\textwidth}
\textbf{Структурная и приведённая формы}

\vspace{0.25em}

\footnotesize

$$
\Gamma(L)y_t = \varepsilon_t
$$

$$
[I - B(L)]y_t = u_t
$$

$$
I - B(L) = \Gamma_0^{-1}\Gamma(L)
$$

$$
\Gamma_0 \Sigma_u \Gamma_0' = \Lambda
$$

\vspace{0.15em}

Здесь $\varepsilon_t$ — структурные шоки, $u_t$ — остатки приведённой VAR, $\Lambda$ — диагональная ковариационная матрица структурных шоков.
\end{column}

\end{columns}

\end{frame}